\section{\label{I_3_C}Le projet \gls{ia} chez Skinsoft} 

\subsection{Un projet global}
La mise en place d'une indexation automatique, visant à faciliter le travail des utilisateurs métier et à le rendre efficace, s'inscrit dans un projet de mise en place de l'intelligence artificielle à une plus large échelle au sein de l'application chez Skinsoft, afin d'automatiser et de faciliter les pratiques de gestion des collections. L'intelligence artificielle (\gls{ia}) peut être définie ainsi : 

\begin{quote} 
	
	Par IA, on entend, dans une acception large et ouverte, la tentative de construire des machines et des systèmes capables de penser et d'agir rationnellement \footcite{zotero-item-584}. 
	
\end{quote} 



L'objectif serait alors, pour un éditeur comme Skinsoft, celui de rendre ces systèmes capables de pouvoir raisonner afin d'effectuer les mêmes actions que l'homme. Dans le cadre d'un système de gestion des collections, cela signifie que les systèmes d'intelligence artificielle devront effectuer des tâches documentaires complexes telles que nous les avons décrites. Cette implémentation doit donc s'inscrire dans cadre défini donnant la direction de développement de ces machines, adapté au contexte patrimonial. 

L'implémentation de l'\gls{ia} au sein de l'application n'est pas seulement issue d'une logique commerciale d'innovation, elle est centrée avant tout sur l'amélioration de l'expérience utilisateur, en vue de rendre l'indexation des collections efficace. L'une des caractéristiques majeures de l'application, nous l'avons évoqué, est sa flexibilité et son adaptabilité aux spécificités des collections et à diverses règles métier. La mise en place de l'\gls{ia} permettrait de mettre cette adaptabilité en place plus efficacement et à plus large échelle. Automatiser la migration de données, le paramétrage de l'application, la navigation dans l'application permettrait à l'entreprise de satisfaire ses clients plus rapidement et de gérer ainsi de multiples contrats. Cette volonté d'implémenter l'\gls{ia} s'inscrit notamment dans l'appel à projets "France 2030" qui couvre la maitrise des technologies numériques et le "déploiement de l'innovation" \footcite{zotero-item-636}. Ce projet finance des initiatives à l'intersection de la culture, du numérique et de l'\gls{ia}. 

\subsection{Identifier le besoin pour instaurer un cadre éthique}

Toutefois, la question du besoin est centrale, pour une utilisation de l'\gls{ia} éthique. Comme le mentionne le Ministère de la culture : "il ne s’agit pas d’utiliser l’IA par défaut, mais d’abord d’examiner le besoin à couvrir" \footcite{2025a}. L'objectif est double pour Skinsoft. La mise en place de l'\GLS{ia} dans l'application permettrait d'une part d'autonomiser les institutions sur leur usage du logiciel et d'autre part d'améliorer l'efficacité de l'indexation de leurs collections. Ainsi, les équipes Skinsoft seraient moins sollicitées pour des difficultés liées à l'usage de l'application qui pourraient être réglés par une machine intelligente. Afin de répondre à ces objectifs, Skinsoft prévoie notamment la mise en place d'une recherche en langage naturel, c'est-à-dire de permettre à l'utilisateur d'effectuer des recherches plus facilement en utilisant un langage naturel---son propre langage--- plutôt qu'un langage contrôlé. Cette fonctionnalité est abordée dans la spécification fonctionnelle intégrée à l'annexe A (\ref{annexe_A}). L'objectif est ainsi de rendre l'utilisation globale de l'application plus simple, et d'améliorer l'efficacité de la gestion de collections. 

Toutefois, l'implémentation de l'\gls{ia} pour la création de contenus structurés soulève des enjeux éthiques dès lors que l'objectif de l'indexation est de produire des connaissances pour la transmission auprès du public. De plus, les pratiques documentaires, nous l'avons vu, sont des pratiques normalisées, structurées par un modèle de données afin d'assurer leur intéropérabilité.  

Ainsi, Skinsoft porte des points d'attention au projet. Les moteurs d'intelligence artificielle implémentés devraient notamment faciliter le travail des utilisateurs et être construits comme des outils de proposition, sans réaliser entièrement eux-mêmes la production de connaissance. L'humain est ainsi toujours intégré au sein de la chaîne de travail dans la définition du projet, en étant le réel décisionnaire final de ce que produit l'\gls{ia}.  

\subsection{Le besoin d'une segmentation et d'une indexation automatisées}

Concrètement, Skinsoft souhaite mettre en place différents moteurs d'intelligence artificielle, entraînés sur des données patrimoniales, et dédiés à l'automatisation partielle de tâches métier dans l'application. Parmi ces tâches, il est prévu l'automatisation de l'indexation au sens large au sein de l'application, et notamment appliquée aux archives audiovisuelles, dans S-Movies. Si, pour les images fixes, peut être mise en place une analyse automatique de leur contenu afin de l'indexer, l'indexation automatique des archives audiovisuelles se découpe,en théorie, en plusieurs sous-tâches, propres à la nature du médium cinématographique.  

Ainsi, le moteur d'intelligence artificielle dédié à l'indexation des archives audiovisuelles au sens large, peut être divisé en plusieurs sous-moteurs, tous dédiés à l'analyse de la coordination du son avec le mouvement des images et leur contenu. Concrètement, cela peut désigner les tâches automatiques suivantes : transcription et analyse du son, reconnaissance des locuteurs, segmentation et analyse des images en mouvement, la transcription du texte présent dans l'image. Les données extraites de ces différentes tâches devront être automatiquement inscrites au sein de référentiels afin d'assurer leur interopérabilité. Par exemple, lors de la reconnaissance d'un locuteur, son nom doit pouvoir être relié à sa notice d'autorité dans l'application.  

Toutefois, le projet d'automatiser l'indexation est issu d'une demande spécifique d'une institution cliente du logiciel, gestionnaire du corpus audiovisuel spécifique que nous avons décrit, constitué de \gls{rush}, nécessitant également un processus de segmentation. Or, la mise en place de la segmentation suppose d’en définir la granularité d'analyse au niveau de la syntaxe audiovisuelle, composée d'éléments divers : le \gls{plan}, la \gls{sequence}, la scène, etc. \footcite{stockinger2011}. Dans le cadre de ce corpus, la segmentation manuelle s'effectue déjà au niveau de la \gls{sequence}. Une \gls{sequence}, en termes cinématographiques, peut être définie comme "une unité d'action, un fragment du film qui raconte en plusieurs plans une suite d'événements isolable dans la construction narrative." \footcite{journot2008}. Dans ce contexte, la segmentation désigne donc la division d'images en mouvements en plusieurs \gls{sequence}s. 

L'institution souhaiterait ainsi avoir à sa disposition un outil qui détecterait automatiquement les changements de \gls{sequence}s, pour pouvoir les décrire individuellement de manière fine, évitant ainsi une description trop globale au niveau du \gls{rush}. Comme nous l'avons évoqué, des marqueurs temporels sont déjà présents dans les données de l'institution, au sein des notices Oeuvres, afin de délimiter les \gls{sequence}s pertinentes. L'objectif est alors celui d'automatiser l'apposition de marqueurs temporels sur la vidéo en détectant les changements majeurs d'action, de lieu, d'objets ou de personnages, marquant le passage d'une séquence à une autre. 

À ce besoin s'ajoute celui d'indexer automatiquement les \gls{sequence}s identifiées par des mots-clés, afin de leur donner une valeur sémantique. En effet, les \gls{sequence}s, au sein de la syntaxe cinématographique, désignent une unité à la fois temporelle et sémantique, marquant un changement de forme significatif au sein du flux audiovisuel, mais également un changement de fond et de sens donné à l'image. Ainsi, la segmentation ne suffit pas à les qualifier de \gls{sequence}s, puisque c'est véritablement leur description qui leur redonne leur sens. Cette description doit ainsi pouvoir être automatisée par la production de mots-clés, puis indexée en lien avec les normes métier employées.

Ainsi, le besoin d'automatisation identifié se divise en deux tâches spécifiques : la segmentation du flux audiovisuel en \gls{sequence}s, et leur indexation par mots-clés. L'automatisation de ces tâches permettrait une accélération du traitement des fonds audiovisuels, et à terme, une mise à disposition plus large auprès d'utilisateurs. Si le projet d'automatiser ces tâches est issu d'un besoin spécifique à uen institution, il pourra être généralisé à toute institution qui préserve des collections audiovisuelles. 
